\BourbakiExerciseGroup

\begin{enumerate}
\def\labelenumi{\arabic{enumi})}
\item
  求有限集上所有可能的滤子。
\item
  若无穷集 X 上的滤子 F 的一切集的交是空集，试证 F 比 X 的有限子集的余集所成的滤子细。
\item
  两个滤子 \(\mathfrak{F}\) 与 \(\mathfrak{G}\)（\(X\) 上的）的交是一切 \(M \cup N\) 所成的集，其中 \(M \in \mathfrak{F}\)，\(N \in \mathfrak{G}\)。
\item
  设 X 是无穷集。试证：X 的有限子集的余集所成的滤子是 X 中各项都不同的无穷序列所相应的诸初等滤子的交。
\item
  设 \(\mathfrak{F}\) 与 \(\mathfrak{G}\) 是集 X 上两个滤子。试证：若 \(\mathfrak{F}\) 与 \(\mathfrak{G}\) 在 X 上所有滤子所成的集中有一个上确界，则这个上确界是一切 \(M \cap N\) 所成的集，其中 \(M \in \mathfrak{F}\)，\(N \in \mathfrak{G}\)。
\item
  若把集 X 上的每个滤子对应到集 \(X' = X \cup \{\omega\}\) 上相应的拓扑（第 5 目），则有以下诸命题：
\end{enumerate}

\begin{enumerate}
\def\labelenumi{\alph{enumi})}
\item
  若拓扑 \(\mathfrak{C}\)（\(\mathbf{X}'\) 上的）比一个滤子 \(\mathfrak{F}\)（\(\mathbf{X}\) 上的）所相应的拓扑细，则 \(\mathfrak{C}\) 或者是离散拓扑，或者是一个比 \(\mathfrak{F}\) 细的滤子所相应的拓扑。逆命题如何？
\item
  与 X 上的一个滤子集 \(\Phi\) 中诸滤子相应的拓扑的下确界，是与属于 \(\Phi\) 的诸滤子之交相应的拓扑。
\item
  设 \(\mathfrak{G}\) 是 X 的子集所成的一个集，\(\mathfrak{G}'\) 是形如 \(\mathbf{M} \cup \{\omega\}\)（其中 \(\mathbf{M} \in \mathfrak{G}\)）的 X' 的子集所成的集。用 \(\mathfrak{G}\) 记 \(\mathfrak{G} \cup \mathfrak{P}(\mathbf{X})\)。若 \(\mathfrak{G}\) 生成一个滤子 \(\mathfrak{F}\)，试证 \(\mathfrak{G}\) 是 X' 上与 \(\mathfrak{F}\) 相应的拓扑的一个子基。若 \(\mathfrak{G}\) 不是滤子子基，\(\mathfrak{G}\) 生成的是什么拓扑？
\item
  \(\mathfrak{F}\) 是 \(\mathbf{X}\) 上的超滤子，当且仅当 \(\mathbf{X}'\) 上相应的拓扑是 \(\mathbf{X}\)-极大的（\(\S 3\)，习题 11）。
\end{enumerate}

\begin{enumerate}
\def\labelenumi{\arabic{enumi})}
\setcounter{enumi}{6}
\item
  在拓扑空间 X 中，X 的非空子集 A 的一切点的邻域滤子的交就是 A 在 X 中的邻域滤子。
\item
  \begin{enumerate}
  \def\labelenumii{\alph{enumii})}
  \tightlist
  \item
    设 \(\Phi\) 是集 \(X\) 上一些拓扑所成的集。试证：一点 \(x \in X\) 关于 \(\Phi\) 中诸拓扑的交的邻域滤子，粗于 \(x\) 关于属于 \(\Phi\) 的诸拓扑的邻域滤子的交。
  \end{enumerate}
\end{enumerate}

\begin{enumerate}
\def\labelenumi{\alph{enumi})}
\setcounter{enumi}{1}
\tightlist
\item
  用 \(\mathfrak{G}_1\) 记拓扑 \(\mathfrak{G}_0(\mathbf{Q}^2)\)（§ 2，习题 7），其中集 \(\mathbf{Q}^2\) 按字典序是一个全序集（《集合论》，
\end{enumerate}

第三章，§ 2，第 6 目）。用 \(\mathcal{G}_2\) 记 \(\mathbf{Q}^2\) 上把 \(\mathcal{G}_1\) 借对称映射 \((\xi, \eta) \to (\eta, \xi)\) 搬运过来所得到的拓扑。试证：若 \(\mathcal{G}\) 是拓扑 \(\mathcal{G}_1, \mathcal{G}_2\) 的交，则 \(\mathbf{Q}^2\) 的一点关于 \(\mathcal{G}\) 的邻域滤子严格地粗于该点关于 \(\mathcal{G}_1\) 与 \(\mathcal{G}_2\) 的邻域滤子的交。

\begin{enumerate}
\def\labelenumi{\arabic{enumi})}
\setcounter{enumi}{8}
\tightlist
\item
  \begin{enumerate}
  \def\labelenumii{\alph{enumii})}
  \tightlist
  \item
    试证：比有限个滤子的交细的每个超滤子，至少比其中一个滤子细（用第 4 目的命题 5）。
  \end{enumerate}
\end{enumerate}

\begin{enumerate}
\def\labelenumi{\alph{enumi})}
\setcounter{enumi}{1}
\tightlist
\item
  给出一个超滤子的例子，它比无穷族超滤子的交细，但不等于这族中任何一个超滤子（考虑无穷集 X 上每个都以单独一点所成的集为基的超滤子所成的族）。
\end{enumerate}

\begin{enumerate}
\def\labelenumi{\arabic{enumi})}
\setcounter{enumi}{9}
\item
  试证：一个超滤子的一切集的交至多含一个点；若含一个点，则该超滤子由含这个点的所有集组成（用第 4 目的命题 5）。
\item
  试证：若集 X 的子集 A 不属于 X 上的超滤子 U，则 U 在 A 上的迹是集 \(\mathfrak{P}(A)\)。
\item
  在一个无穷集上，试证与各项都不同的序列相应的初等滤子不是超滤子。
\item
  设 \(f\) 是集 \(\mathbf{X}\) 到集 \(\mathbf{X}'\) 中的一个映射。则 \(f\) 是单射的充要条件是：对每个滤子基 \(\mathfrak{B}\)（\(\mathbf{X}\) 上的），\((f\overline{f}^{1}(\mathfrak{B}))\) 是一个与 \(\mathfrak{B}\) 等价的滤子基。
\item
  试证：若 \(f\) 是集 \(\mathbf{X}\) 到集 \(\mathbf{X}'\) 上的一个映射，则任一滤子的像 \(f(\mathfrak{F})\)（\(\mathfrak{F}\) 是 \(\mathbf{X}\) 上的滤子）是 \(\mathbf{X}'\) 上的一个滤子。
\end{enumerate}

¶ 15) 设 \(n \to f(n)\) 是 \(\mathbf{N}\) 到 \(\mathbf{N}\) 上的一个映射，使得 \(\overline{f}(m)\)（对每个 \(m \in \mathbf{N}\)）是有限的。若 \((x_n)\) 是集 \(\mathbf{X}\) 中的任一序列，令 \(y_n = x_{f(n)}\)。试证与序列 \((x_n)\) 和 \((y_n)\) 相应的初等滤子是相同的。

由此推出：若 \((a_{n})\) 与 \((b_{n})\) 是集 X 中元素的两个序列，且与 \((b_{n})\) 相应的滤子比与 \((a_{n})\) 相应的滤子细，则与 \((b_{n})\) 相应的滤子与 \((a_{n})\) 的某个子列所相应的滤子相同。

¶ 16) 设 \(\Phi\) 是集 X 上初等滤子所成的一个可数全序集。试证存在一个初等滤子，它比 \(\Phi\) 中所有滤子都细（试证 \(\Phi\) 中诸滤子的并有一个可数基）。

\begin{enumerate}
\def\labelenumi{\arabic{enumi})}
\setcounter{enumi}{16}
\tightlist
\item
  设 X 是一个格（《集合论》，第三章，§ 1，第 13 目）。X 的非空子集 F 称为预滤子，如果它满足下列条件：(i) 若 \(x \in \mathbf{F}\) 且 \(y \geqslant x\)，则 \(y \in \mathbf{F}\)；(ii) 若 \(x \in \mathbf{F}\) 且 \(y \in F\)，则 \(\inf (x, y) \in F\)；(iii) \(F \neq X\)。于是，按包含关系排序时，集 Y 的子集全体上的预滤子恰好就是 Y 上的滤子。预滤子 F 称为素的，如果关系 \(\sup(x, y) \in F\) 蕴含 \(x \in F\) 或 \(y \in F\)。关于与 X 的序相反的序的预滤子称为 X 上的余预滤子。
\end{enumerate}

\begin{enumerate}
\def\labelenumi{\alph{enumi})}
\item
  X 中不独由最小元组成的子集 A 的上界所成的集是一个预滤子。给出一些不是上界集的预滤子的例子。
\item
  试证：若 F 是素预滤子，则 CF 是素余预滤子。
\item
  试证：若 X 有最小元 0，则每个预滤子都含于一个极大预滤子（用 Zorn 引理，并注意 0 不属于任何预滤子）。
\item
  求全序集 X 上所有的预滤子，并证它们都是素的。由此给出一些不是极大的素预滤子的例子。
\item
  设 X 是由 5 个元素组成的有序集，记作 0, 1, a, b, c，其序关系为 0 \textless{} a \textless{} 1、0 \textless{} b \textless{} 1、0 \textless{} c \textless{} 1。试证 X 是一个格，且除 \(\{1\}\) 外的预滤子都是极大的但不是素的。
\item
  设 X 是有最小元 0 的格。试证：若 \(F \subset X\) 是一个预滤子，而 a 是 X 的一个不属于 F 的元素，则存在包含 F 与 a 的预滤子的充要条件是 \(\inf (a, x) \neq 0\)（对一切 \(x \in F\)）；这时包含 F 与 a 的最小预滤子是一切满足下列条件的 \(y \in X\) 所成的集：\(y \geqslant \inf (a, x)\)（对 F 的至少一个元素 x）。
\end{enumerate}

¶ 18) 设 X 是一个分配格（《集合论》，第三章，§ 1，习题 16），它有一个最小元 0。

\begin{enumerate}
\def\labelenumi{\alph{enumi})}
\item
  若 \(a\) 是 \(\mathbf{X}\) 的一个既约元（《集合论》，第三章，§ 4，习题 7），试证 \(a\) 的一切上界所成的集是一个素预滤子。给出一个素预滤子的例子，它不是 \(\mathbf{X}\) 的既约元的上界集{[}参见习题 17 d){]}。
\item
  试证 X 上每个极大预滤子都是素的{[}用习题 17f){]}。给出分配格上的素而非极大的预滤子的一些例子{[}习题 17d){]}。
\item
  试证每个预滤子 F 都是包含 F 的素预滤子的交{[}若 \(a \notin F\)，试证包含 F 但不含 \(a\) 的预滤子所成的集中的一个极大元 U 是素的，这要用习题 17f){]}。
\item
  设 \(\Omega\) 是 X 上素预滤子所成的集。对每个 \(x \in \mathbf{X}\)，使对应于子集 \(\mathbf{A}_x\)（\(\Omega\) 中由满足 \(x \in \mathbf{F}\) 的素预滤子 F 组成的子集）。试证映射 \(x \to \mathbf{A}_x\)（X 到 \(\mathfrak{P}(\Omega)\) 中的映射）是单射且保序，\(\mathbf{A}_{\mathrm{inf}(x,y)} = \mathbf{A}_x \cap \mathbf{A}_y\)，并且
\end{enumerate}

\[
\mathbf {A} _ {\sup (x, y)} = \mathbf {A} _ {x} \cup \mathbf {A} _ {y}.
\]

\begin{enumerate}
\def\labelenumi{\arabic{enumi})}
\setcounter{enumi}{18}
\item
  设 X 是有最小元的格。试证：若 X 上每个预滤子都是包含它的素预滤子的交，则 X 是分配的（若映射 \(x \rightarrow A_{x}\) 像习题 18 d) 中那样定义，则这个映射仍具有那里所述的性质）。
\item
  格 X 称为布尔代数，如果它是分配的，有最小元 0 和最大元 1，并且对每个 \(x \in \mathbf{X}\)，存在元素 \(x' \in \mathbf{X}\) 使 \(\inf (x, x') = 0\) 且 \(\sup (x, x') = 1\)；这样的元素称为 \(x\) 在 \(\mathbf{X}\) 中的补元。
\end{enumerate}

\begin{enumerate}
\def\labelenumi{\alph{enumi})}
\item
  设 \(x \to \mathbf{A}_x\) 是布尔代数 \(\mathbf{X}\) 到集 \(\Omega\)（\(\mathbf{X}\) 上素预滤子所成的集）的子集全体中的单射，如习题 18 d) 中所定义。试证：若 \(x'\) 是 \(x\) 在 \(\mathbf{X}\) 中的一个补元，则 \(\mathbf{A}_{x'} = \complement \mathbf{A}_x\)（在 \(\Omega\) 中）；由此推出：元素 \(x\) 只有一个补元 \(x'\)，\(x'\) 的补元是 \(x\)，而且 \(\inf (x, y)\) 的补元是 \(\sup (x', y')\)；对这些命题给出直接的证明。布尔环 \(\mathbf{A}\)（相应于 \(\mathbf{X}\) 的环）的理想，除 \(\mathbf{A}\) 本身外，就是 \(\mathbf{X}\) 上的余预滤子。
\item
  试证：在布尔代数 \(\mathbf{X}\) 中，每个素预滤子都是极大的（注意，对每个 \(x \in \mathbf{X}\)，每个素预滤子必含 \(x\) 或其补元）。
\item
  反之，设 X 是一个分配格，它有最小元 0 和最大元 1，并且 X 上每个素预滤子都是极大的。试证 X 是一个布尔代数。（对元素 \(x \in X\)，考虑由那些 \(y \in X\)（满足 \(\inf(x, y) = 0\)）所组成的余预滤子 C（习题 17）；若 \(x\) 没有补元，则有一个包含 \(x\) 与 C 的极大余预滤子 M；注意 M 在 X 中的余 U 是一个素预滤子{[}习题 17 b) 与 18 b){]}，并用习题 17 f)）。d) 试证：若 X 是拓扑空间，则 X 中既开又闭的子集所成的集 B 按包含关系排序是一个布尔代数。若取 X 为与 Fréchet 滤子相应的拓扑空间（第 5 目），试证这样得到的布尔代数 B 是可数无穷的，因此不能同构于一个集 A 的所有子集所成的有序集 P(A)。
\end{enumerate}
% semantic-fragments: legacy-input-seam
\relax{}
